\documentclass[11pt,a4paper]{article}
\usepackage[margin=0.72in,top=0.70in,bottom=0.70in]{geometry}
\usepackage{amsmath,amssymb}
\usepackage{times}
\usepackage{enumitem}
\usepackage{fancyhdr}

\pagestyle{fancy}
\fancyhf{}
\fancyhead[R]{\small\thepage}
\renewcommand{\headrulewidth}{0pt}
\setlength{\headheight}{14.5pt}

\setcounter{page}{6}

\setlength{\parindent}{0pt}
\setlength{\parskip}{0.35em}

% Compact display equations
\setlength{\abovedisplayskip}{4pt}
\setlength{\belowdisplayskip}{4pt}

\begin{document}

% =========================================================================
% PAGE 1 (Matching Screenshot 1 / Page 22)
% =========================================================================

\subsection*{4.5\quad Algorithm of PrismSpace Multi-Agent System}

\textbf{Input:}
\begin{itemize}[leftmargin=1.5em, itemsep=0.25em, topsep=0.10em]
    \item Natural language user objective, task specifications, and operational constraints ($\mathcal{T}_{\text{user}}$).
    \item Multi-source curated corpora and execution traces ($\mathcal{D}$), spanning LMSYS Arena, RouterBench, WildGuard, CMU Agent Trajectories, and PrismSpace Runtime Events.
    \item Curated preference alignment pairs ($\mathcal{D}_{\text{pref}} = \{(x, y_w, y_l)\}$), containing user prompt $x$, winning/chosen agent completion $y_w$, and losing/rejected completion $y_l$.
    \item Global multi-agent swarm and tool registry ($\mathcal{G}_{\text{global}} = (V, E)$), comprising the Supervisor Hub, specialized worker agents (Coding, Research, Validator, Planner), and Model Context Protocol (\textbf{MCP}: open standard defining structured JSON-RPC communication between agents and tools) endpoints.
\end{itemize}

\textbf{Output:}
\begin{itemize}[leftmargin=1.5em, itemsep=0.25em, topsep=0.10em]
    \item Predicted intent category ($\hat{y}_{\text{intent}}$) and optimal Large Language Model (\textbf{LLM}: deep neural network trained for language comprehension and generation) provider $\hat{y}_{\text{provider}}$ (e.g., Groq, NVIDIA, OpenAI, Anthropic, Google).
    \item Swarm agent dispatch vector ($\hat{\mathbf{y}}_{\text{agents}} \in \{0, 1\}^{|V|}$) and safety approval gate decision ($\hat{y}_{\text{approval}} \in \{0, 1\}$).
    \item Active collaboration topology Directed Acyclic Graph ($\mathcal{G}_{\text{active}}$ [\textbf{DAG}: a directed graph without closed cycles, establishing execution dependencies]) with cost ($\hat{c}$) and latency ($\hat{\ell}$) predictions.
    \item Preference-aligned adapter weights ($\Delta \mathbf{W}$) optimizing response quality, policy compliance, and safety margin.
    \item Performance evaluation metrics: \textbf{Accuracy} (correct prediction ratio), \textbf{Macro-F1} (unweighted harmonic mean of precision and recall across all classes), \textbf{ROC-AUC} (Area Under the Receiver Operating Characteristic Curve measuring discrimination across thresholds), \textbf{ECE} (Expected Calibration Error: average gap between confidence and accuracy), \textbf{Brier Score} (probability forecast error), \textbf{$R^2$} (Coefficient of Determination: explained variance), and \textbf{MAE} (Mean Absolute Error: average magnitude of prediction errors).
\end{itemize}

\vspace{0.6em}
\textbf{Step 1: Data Acquisition and Preprocessing}
\begin{enumerate}[leftmargin=1.5em, itemsep=0.28em, topsep=0.12em]
    \item Ingest heterogeneous records from multi-source corpora ($\mathcal{D}$) covering intent traces, safety pairs, tool telemetry, and conversational benchmarks.
    \item Clean and normalize feature values (strip corrupt 8-bit Unicode Transformation Format [\textbf{UTF-8}: standard variable-width character encoding] sequences, normalize whitespace, remove duplicates).
    \item Encode categorical agent attributes and tool capabilities (e.g., persona role, execution scope).
    \item Standardize all numerical features to zero mean and unit variance:
    \[
    z_i = \frac{x_i - \mu_x}{\sigma_x}
    \]
    \item Split the dataset into 80\% training ($\mathcal{D}_{\text{train}}$), 10\% validation ($\mathcal{D}_{\text{val}}$), and 10\% testing ($\mathcal{D}_{\text{test}}$) partitions.
\end{enumerate}

\newpage

% =========================================================================
% PAGE 2: Graph Construction & Attention-Based Message Passing
% =========================================================================

\textbf{Step 2: Graph Construction}
\begin{enumerate}[leftmargin=1.5em, itemsep=0.30em, topsep=0.12em]
    \item Treat each agent persona and MCP tool $v_i$ as a node in the swarm graph ($V$).
    \item Compute pairwise similarity between feature vectors using \textbf{Pearson correlation}:
    \[
    s_{ij} = \frac{\mathrm{Cov}(\mathbf{x}_i, \mathbf{x}_j)}{\sigma_{\mathbf{x}_i}\sigma_{\mathbf{x}_j}}
    \]
    where $\mathrm{Cov}$ denotes Covariance (statistical measure of joint directional variance between two feature distributions).
    \item Define adjacency matrix $\mathbf{A}$ such that:
    \[
    A_{ij} = \begin{cases} 1, & \text{if } s_{ij} \ge \tau \\ 0, & \text{otherwise} \end{cases}
    \]
    where $\tau$ is the similarity threshold.
    \item Construct the graph $\mathcal{G} = (V, E)$ with edges $E$ derived from adjacency matrix $\mathbf{A}$.
\end{enumerate}

\vspace{0.6em}
\textbf{Step 3: Graph Attention Network (\textbf{GAT}) Initialization}
\begin{enumerate}[leftmargin=1.5em, itemsep=0.30em, topsep=0.12em]
    \item Initialize node embeddings $\mathbf{h}_i = \mathbf{x}_i$ for all nodes $v_i \in V$ in the Graph Attention Network (\textbf{GAT}: neural architecture utilizing self-attention to weight graph neighbor connections).
    \item Initialize weight matrix $\mathbf{W}$ and attention vector $\mathbf{a}$ with random values.
    \item Set Learning Rate (\textbf{LR}: optimization step size governing parameter updates), number of attention heads ($K$), and dropout probability ($p$).
\end{enumerate}

\vspace{0.6em}
\textbf{Step 4: Attention-Based Message Passing}
\begin{enumerate}[leftmargin=1.5em, itemsep=0.30em, topsep=0.12em]
    \item For each node $v_i$, compute attention scores for all its neighbors $v_j \in \mathcal{N}(i)$:
    \[
    e_{ij} = \mathrm{LeakyReLU}\left(\mathbf{a}^T [\mathbf{W} \mathbf{h}_i \parallel \mathbf{W} \mathbf{h}_j]\right)
    \]
    where $\mathrm{LeakyReLU}$ denotes Leaky Rectified Linear Unit (activation function with a small non-zero slope for negative values to prevent dead neurons), and $\parallel$ represents feature concatenation.
    \item Normalize attention coefficients using the softmax function:
    \[
    \alpha_{ij} = \frac{\exp(e_{ij})}{\sum_{k \in \mathcal{N}(i)} \exp(e_{ik})}
    \]
    \item Update the hidden representation of node $v_i$:
    \[
    \mathbf{h}_i' = \sigma\left(\sum_{j \in \mathcal{N}(i)} \alpha_{ij} \mathbf{W} \mathbf{h}_j\right)
    \]
    \item Apply dropout regularization and repeat for all layers.
\end{enumerate}

\newpage

% =========================================================================
% PAGE 3: Classification & Preference Alignment (ORPO)
% =========================================================================

\textbf{Step 5: Classification and Loss Computation}
\begin{enumerate}[leftmargin=1.5em, itemsep=0.30em, topsep=0.12em]
    \item Feed the final embeddings $\mathbf{h}_i'$ into a fully connected classification layer:
    \[
    \hat{y}_i = \mathrm{Sigmoid}(\mathbf{W}_c \mathbf{h}_i')
    \]
    \item Compute the Binary Cross-Entropy (\textbf{BCE}: loss function quantifying divergence between true binary labels and predicted probabilities) loss:
    \[
    L = -\frac{1}{n} \sum_{i=1}^n \left[ y_i \log(\hat{y}_i) + (1 - y_i) \log(1 - \hat{y}_i) \right]
    \]
    \item Backpropagate the loss through the network and update parameters $\theta = \{\mathbf{W}, \mathbf{a}, \mathbf{W}_c\}$ using the Adam optimizer (\textbf{Adam}: Adaptive Moment Estimation, an algorithm combining momentum and RMSProp gradient scaling).
\end{enumerate}

\vspace{0.6em}
\textbf{Step 6: Preference Alignment and Reward Optimization (\textbf{ORPO})}
\begin{enumerate}[leftmargin=1.5em, itemsep=0.28em, topsep=0.12em]
    \item Ingest curated preference triples $\mathcal{D}_{\text{pref}} = \{(x, y_w, y_l)\}$ comprising user prompt $x$, winning/chosen agent completion $y_w$, and losing/rejected completion $y_l$.
    \item Parameterize the base causal language model via Parameter-Efficient Fine-Tuning (\textbf{PEFT}: techniques adapting foundation models with minimal trainable weights) using Low-Rank Adaptation (\textbf{LoRA}: parameter-efficient fine-tuning decomposing weight updates into low-rank matrices $\Delta \mathbf{W} = \frac{\alpha}{r} \mathbf{B}\mathbf{A}$).
    \item Compute the Supervised Fine-Tuning (\textbf{SFT}: training language models via next-token teacher forcing) Negative Log-Likelihood (\textbf{NLL}: loss measuring prediction error of ground-truth token sequences) on the chosen response $y_w$:
    \[
    \mathcal{L}_{\mathrm{SFT}}(\theta) = -\frac{1}{|y_w|} \sum_{t=1}^{|y_w|} \log P_\theta(y_{w,t} \mid x, y_{w,<t})
    \]
    \item Formulate the length-normalized sequence probability $P_\theta(y \mid x) = \left(\prod_{t=1}^{|y|} P_\theta(y_t \mid x, y_{<t})\right)^{\frac{1}{|y|}}$ and define the generative odds of sequence $y$ given prompt $x$:
    \[
    \mathrm{odds}_\theta(y \mid x) = \frac{P_\theta(y \mid x)}{1 - P_\theta(y \mid x)}
    \]
    \item Compute the Odds Ratio preference loss to penalize disfavored agent behaviors without requiring a reference model:
    \[
    \mathcal{L}_{\mathrm{OR}}(\theta) = -\log \sigma \left( \log \frac{\mathrm{odds}_\theta(y_w \mid x)}{\mathrm{odds}_\theta(y_l \mid x)} \right)
    \]
    where $\sigma$ denotes the Sigmoid activation function.
    \item Optimize the unified Odds Ratio Preference Optimization (\textbf{ORPO}: reference-model-free alignment technique directly penalizing disfavored generation odds during supervised fine-tuning) objective:
    \[
    \mathcal{L}_{\mathrm{ORPO}}(\theta) = \mathcal{L}_{\mathrm{SFT}}(\theta) + \beta \cdot \mathcal{L}_{\mathrm{OR}}(\theta)
    \]
    where $\beta$ is the preference penalty scaling hyperparameter (e.g., $\beta = 0.1$).
\end{enumerate}

\newpage

% =========================================================================
% PAGE 4: Runtime Inference, Swarm Execution, and UI Generation
% =========================================================================

\textbf{Step 7: Real-Time Runtime Inference, Swarm Execution, and UI Generation}
\begin{enumerate}[leftmargin=1.5em, itemsep=0.30em, topsep=0.12em]
    \item \textbf{User Input Ingestion}: Capture natural language task objective $\mathcal{T}_{\text{user}}$ submitted via the Next.js User Interface (\textbf{UI}: client-facing presentation layer for reactive interaction) and transmit via Hypertext Transfer Protocol (\textbf{HTTP}: standard web communication protocol) POST to the FastAPI Hive Bridge Application Programming Interface (\textbf{API}: formal programmatic contract enabling service communication).
    \item \textbf{Query Vectorization}: Extract lexical, semantic, and telemetry features:
    \[
    \mathbf{x}_{\text{query}} = [\mathbf{v}_{\text{TF-IDF}}(\mathcal{T}_{\text{user}}) \parallel \mathbf{e}_{\text{dense}}(\mathcal{T}_{\text{user}}) \parallel \mathbf{f}_{\text{meta}}]
    \]
    where $\parallel$ denotes vector concatenation, and $\mathbf{f}_{\text{meta}}$ encodes query length, token count, and domain complexity.
    \item \textbf{GAT Model Routing Inference}: Compute hidden query state $\mathbf{h}_{\text{query}}'$ to predict intent, optimal LLM provider, and active agent dispatch mask:
    \[
    \hat{y}_{\text{intent}}, \; \hat{y}_{\text{provider}} = \mathrm{argmax}\left(\mathrm{Softmax}(\mathbf{W}_{\text{route}} \mathbf{h}_{\text{query}}')\right), \quad \hat{\mathbf{y}}_{\text{agents}} = \mathbb{I}\left(\mathrm{Sigmoid}(\mathbf{W}_{\text{agent}} \mathbf{h}_{\text{query}}') \ge \tau_{\text{dispatch}}\right)
    \]
    where $\mathbb{I}(\cdot)$ is the indicator function and $\tau_{\text{dispatch}}$ is the routing activation threshold.
    \item \textbf{Active DAG Task Scheduling}: Synthesize Directed Acyclic Graph $\mathcal{G}_{\text{active}} = (V_{\text{active}}, E_{\text{active}})$ over active agents $V_{\text{active}} = \{v_i \in V \mid \hat{\mathbf{y}}_{\text{agents}, i} = 1\}$, sorted into execution sequence $\pi = \mathrm{TopologicalSort}(\mathcal{G}_{\text{active}})$.
    \item \textbf{Autonomous ReAct Swarm Execution}: For each agent $v_i \in \pi$, execute Reasoning and Acting (\textbf{ReAct}: paradigm interleaving chain-of-thought reasoning with tool execution actions) invoking Model Context Protocol (\textbf{MCP}) endpoints via JavaScript Object Notation Remote Procedure Call (\textbf{JSON-RPC}: lightweight remote procedure call protocol):
    \[
    a_{i, t} = \mathrm{LLM}_{\hat{y}_{\text{provider}}}(\mathcal{T}_{\text{user}}, h_{i, <t}), \quad o_{i, t} = \mathrm{Exec}_{\text{MCP}}(a_{i, t})
    \]
    where $a_{i,t}$ represents the action/tool call payload and $o_{i,t}$ denotes the environment observation.
    \item \textbf{ORPO Safety Gate Verification}: Evaluate candidate completion $\hat{y}_{\text{candidate}}$ using the preference policy:
    \[
    \mathrm{Score}_{\text{policy}} = \log \frac{\mathrm{odds}_\theta(\hat{y}_{\text{candidate}} \mid \mathcal{T}_{\text{user}})}{\mathrm{odds}_\theta(y_{\text{ref}} \mid \mathcal{T}_{\text{user}})}
    \]
    Approve output if $\mathrm{Score}_{\text{policy}} \ge \tau_{\text{safety}}$; otherwise reroute to Validator agent for policy correction.
    \item \textbf{Real-Time Streaming to UI and Client Persistence}: Stream output tokens and agent thought states to the client via Server-Sent Events (\textbf{SSE}: persistent unidirectional streaming protocol):
    \[
    \text{State}_{\text{UI}}^{(t)} = \text{State}_{\text{UI}}^{(t-1)} \cup \{\Delta_{\text{token}}^{(t)}, \text{Card}_{\text{agent}}\}
    \]
    Upon stream completion, commit conversation turns, generated artifacts, and telemetry to the browser's local IndexedDB (Dexie.js) for reactive rendering via \texttt{useLiveQuery}.
\end{enumerate}

\newpage

% =========================================================================
% PAGE 5: Evaluation, Visualization, and Summary
% =========================================================================

\textbf{Step 8: Model Evaluation}
\begin{enumerate}[leftmargin=1.5em, itemsep=0.28em, topsep=0.12em]
    \item Evaluate the trained routing and alignment models on validation and test partitions.
    \item Compute classification and preference evaluation metrics:
    \begin{itemize}[leftmargin=1.2em, itemsep=0.22em, topsep=0.10em]
        \item \textbf{Accuracy}: Ratio of correct predictions over total predictions:
        \[
        \mathrm{Accuracy} = \frac{TP + TN}{TP + TN + FP + FN}
        \]
        where:
        \begin{itemize}[leftmargin=1.0em, itemsep=0.08em, topsep=0.04em]
            \item \textbf{TP}: True Positives (positive cases correctly identified as positive).
            \item \textbf{TN}: True Negatives (negative cases correctly identified as negative).
            \item \textbf{FP}: False Positives (negative cases incorrectly classified as positive; Type I error).
            \item \textbf{FN}: False Negatives (positive cases incorrectly classified as negative; Type II error).
        \end{itemize}
        \item \textbf{AUC}: Area Under the Curve of the Receiver Operating Characteristic (\textbf{ROC}: curve plotting True Positive Rate vs. False Positive Rate across thresholds) to assess discrimination capability.
        \item \textbf{AUPRC}: Area Under the Precision-Recall Curve (metric evaluating precision-recall trade-offs especially under severe class imbalance).
    \end{itemize}
    \item Save the trained model parameters, LoRA adapters, and generate prediction outputs.
\end{enumerate}

\vspace{0.6em}
\textbf{Step 9: Visualization and Interpretation}
\begin{itemize}[leftmargin=1.5em, itemsep=0.28em, topsep=0.10em]
    \item Visualize learned attention weights to identify key routing features contributing to optimal swarm delegation.
    \item Generate performance plots (loss curve, ROC, precision-recall, reward margin distribution).
    \item Analyze high-attention edges to reveal biologically meaningful and task-optimal variant relationships.
\end{itemize}

\vspace{0.8em}
\textbf{Algorithm Summary}

The algorithm coordinates intelligent model routing, preference alignment, and real-time execution into a cohesive user workflow. By modeling specialized agents and tools through a Graph Attention Network (\textbf{GAT}), the system learns how agent capabilities complement one another, dynamically routing each incoming request to the most capable and cost-effective language model provider. To guarantee that generated actions remain safe and compliant with user intentions, Odds Ratio Preference Optimization (\textbf{ORPO}) directly refines model behaviors against curated preference pairs without reference model overhead. Once an objective is submitted in the web interface, the algorithm deconstructs it into a structured plan, orchestrates autonomous agents to execute tools, verifies safety policies, and streams responses incrementally into the browser's local IndexedDB (Dexie.js). This design ensures transparent, dependable, and responsive multi-agent task execution across both research workflows and enterprise applications.

\end{document}
